\bookchapter{Chapter 36: Phones and Smart Home — Make Your House Think}

\chapterintro{Your phone can do one more increasingly important job: be the remote for your whole home. From a light bulb or plug to a door lock, smart home gear is moving from “geek toy” to everyday setup. This chapter sorts out protocols, ecosystems, hubs, and starter kits, then closes with a safety checklist so “smart” doesn't backfire.}

\booksection{Smart home protocols: Pick the right one}

Before buying a smart device, check which protocols it supports — that decides whether your phone and ecosystem can control it. Four mainstream options and their personalities:

\begin{itemize}
  \item Matter: Released in 2022 by the CSA (Connectivity Standards Alliance), backed by Apple, Google, Amazon, Samsung, and others. One goal: make every brand speak the same language. Rapidly becoming the unified standard from 2023–2026. Matter devices can join Apple HomeKit, Google Home, Amazon Alexa, and Mi Home at once — no more single-ecosystem lock-in.
  \item Thread: Matter's underlying network "plumbing." Low power, self-healing mesh (if a node drops, signals reroute), fast and smooth. Devices with a Thread border router (HomePod mini, Apple TV 4K, Xiaomi central hub) feel noticeably more stable.
  \item Wi-Fi: Big bandwidth, connects directly to the router, but power-hungry and can clog channels with many devices. Good for cameras and robot vacuums that move lots of data.
  \item Bluetooth (BLE): Easy setup, power-efficient, but short range and no mesh. Good for locks and temp/humidity sensors with tiny payloads.
  \item Zigbee: Classic low-power mesh protocol; needs a dedicated hub. Devices are cheap and stable, but cross-brand interop is weak. Gradually replaced by Matter + Thread.
\end{itemize}

\begin{tcbinfo}{Rule of thumb}
Buying smart gear in 2026? Prefer the “Matter over Thread” label — the most future-proof choice today. Older Wi-Fi or Bluetooth devices don't need to be tossed; keep using them and move to Matter on new purchases.

\end{tcbinfo}

\booksection{The big three ecosystems}

Protocols are the “language”; ecosystems are the “club.” Even with Matter, you'll still pick one main ecosystem for automations and voice. Quick comparison:

\begin{center}
\begin{tabularx}{\linewidth}{|X|X|X|X|}
\hline
\textbf{Ecosystem} & \textbf{Strengths} & \textbf{Weaknesses} & \textbf{Best for} \tabularnewline
\hline
Apple HomeKit & Strongest privacy (mostly local), smooth iPhone/Siri/Shortcuts integration & Fewer, pricier devices; needs HomePod or Apple TV as hub & Apple households that value privacy and polish \tabularnewline
Mi Home (Xiaomi) & Widest device range, best value, solid Xiao AI & More closed ecosystem, more cloud data, more ads & Budget buyers who want one-stop whole-home setup \tabularnewline
Huawei HiLink (HarmonyOS) & Seamless with HarmonyOS phone/tablet/watch; device range growing fast & Fewer overseas devices; ecosystem still catching up & Huawei phone users who want Super Device multi-device flow \tabularnewline
\hline
\end{tabularx}
\end{center}

\begin{tcbinfo}{How to choose}
Follow your phone: iPhone HomeKit first. Huawei HiLink first. Other Android Mi Home is easiest. With Matter, devices from different ecosystems can still be controlled by the same phone — no need to pick sides all at once.

\end{tcbinfo}

\booksection{Your phone as the home hub}

In most starter setups, the phone is your “home hub” — control lights, AC, curtains, and security from an app, and build automations so the house “moves on its own”:

\begin{itemize}
  \item Device control: Open HomeKit/Mi Home/Smart Life to toggle lights, change color temperature, set AC temperature, open/close curtains, and view cameras — locally or remotely.
  \item Automation scenes: Combine "condition action." Three favorites:
  \item Away mode: Last person leaves (location leaves or door locks) lights off + AC off + plugs off + cameras on + arm security.
  \item Sleep mode: Every night 23:00 main lights off + curtains close + AC to 26°C + night light on.
\end{itemize}

NFC tap: Stick an NFC tag by the door; tap the phone to trigger a scene (tap once for entry lights, twice for sleep). On iPhone use Shortcuts automation; on Android use Wallet or the brand app to write the tag.

\begin{tcbpractice}{Try it: Build a "Coming home" scene}
In your smart home app's “Automation,” create one: trigger “arrive at place” (home), actions “entry light on,” “AC to 26°C,” “curtains close.” Save it, go out for a walk, and see if the lights turn on when you return — that first success feels better than ordering new gear.

\end{tcbpractice}

\booksection{Starter kit picks}

You don't need everything at once. These four cover about 80\% of daily needs:

\begin{itemize}
  \item Smart speaker: Xiao AI (Mi Home, great value) / HomePod mini (Apple, also a Thread border router) / Alexa (broad overseas ecosystem). Note: The speaker picks your voice ecosystem — decide the ecosystem before buying.
  \item Smart bulbs: Yeelight (Mi Home/Matter, adjustable color temperature and brightness, cheap) / Philips Hue (great color and stability, but needs a Hue hub and costs more). Note: Check socket type (E27/B22) and whether you need color changing.
  \item Smart plugs: Xiaomi Smart Plug (cheap, energy stats) / TP-Link Kasa (good Matter support overseas). Note: For high-power appliances (water heater, space heater), always check the rated wattage and don't overload.
  \item Smart locks: Loock/Xiaomi (value and Mi Home integration) / Kaadas (traditional lock maker, solid mechanics). Note: Check Matter/your ecosystem support, battery life, and where the emergency keyhole is.
\end{itemize}

\begin{tcbinfo}{Suggested buying order}
Smart speaker first (locks your ecosystem) then bulbs and plugs (low barrier, instant feedback) finally locks and cameras (security-related, choose carefully). Add one category at a time so you don't end up with gear you can't use.

\end{tcbinfo}

\booksection{Safety basics}

More smart devices mean a bigger attack surface. These four lines are non-negotiable — check them one by one:

\begin{enumerate}
  \item Change default passwords: Factory passwords for cameras, routers, and locks are public online. First thing after unboxing: set a long, unique password.
  \item Update firmware promptly: Vendor updates often fix security holes. Turn on auto-update in the app, or check manually at least monthly.
  \item Network isolation: Put IoT devices on a separate Wi-Fi (or guest network) away from your phone and computer. If one device is compromised, it's harder to pivot to your important data.
  \item Privacy: Enable camera "privacy zones" (skip bedrooms and bathrooms), turn off 24/7 cloud recording, and don't leave home camera feeds open to third-party apps long-term.
\end{enumerate}

\textbf{Don't step in these traps}\par

\begin{itemize}
  \item Don't use remote unlock for strangers — prefer temporary passwords that expire.
  \item Don't point cameras at beds or out the window — that's your own privacy leaking.
  \item Don't buy no-name "white label" gear for a bargain — it often never gets firmware updates, basically a back door into your home.
\end{itemize}

\begin{tcbsummary}{Chapter recap}
\begin{itemize}
  \item Matter + Thread is the 2023–2026 direction — look for that label on new gear
  \item Each ecosystem has a strength: Apple for privacy/experience, Mi Home for value, Huawei for multi-device flow
  \item Your phone can be the hub — set up Coming home / Away / Sleep and the house runs itself
  \item Starter four: smart speaker, bulbs, plugs, locks — add them in that order
  \item Change default passwords, update firmware, isolate the network, privacy-mask cameras — four safety lines you can't skip
\end{itemize}

\end{tcbsummary}
